% Generated draft: requires mathematical and source review.
\documentclass[11pt]{article}
\usepackage[margin=25mm]{geometry}
\usepackage{fontspec}
\setmainfont{Latin Modern Roman}
\usepackage{amsmath,amssymb,mathtools,unicode-math}
\setmathfont{Latin Modern Math}
\usepackage{enumitem,needspace,xurl,hyperref}
\setlength{\parindent}{0pt}
\setlength{\parskip}{5pt}
\setlength{\emergencystretch}{3em}
\newcommand{\N}{\mathbb{N}}
\newcommand{\Z}{\mathbb{Z}}
\newcommand{\Q}{\mathbb{Q}}
\newcommand{\R}{\mathbb{R}}
\newcommand{\C}{\mathbb{C}}
\newcommand{\F}{\mathbb{F}}
\newcommand{\A}{\mathbb{A}}
\newcommand{\PP}{\mathbb{P}}
\newcommand{\E}{\mathbb{E}}
\newcommand{\eps}{\varepsilon}
\newcommand{\dd}{\,\mathrm{d}}
\begin{document}
\section*{UnsolvedMath ID: 20002836}
\textbf{Source problem code:} catalogue entry
\subsection{Definitions and mathematical statement}
\textbf{Source problem and recovered antecedent.} In Section 8, ``Kisin: Pseudo-representations'', recast dimension-\(d\) trace pseudocharacter theory without assuming \((d!)^{-1}\in A\). The source warns that this is not merely a question about divided powers. The preceding setup, supplied by the verification, has a group \(G\), commutative coefficient ring \(A\), or more generally a unital \(A\)-algebra \(R\). The trailing ``8.2'' is a neighboring section marker, not an additional question.

\textbf{Reported replacement by determinant laws.} For arbitrary commutative unital \(A\), a degree-\(d\) determinant is a multiplicative homogeneous polynomial law \(D:R\to A\). Concretely it is a family of maps
\[
 D_B:R\otimes_A B\to B
\]
natural in every commutative \(A\)-algebra \(B\), with \(D_B(1)=1\), \(D_B(xy)=D_B(x)D_B(y)\), and \(D_B(bx)=b^dD_B(x)\), in the polynomial-law sense. The characteristic coefficients are specified by
\[
 D_{A[t]}(t1-r)=\sum_{i=0}^d(-1)^i\Lambda_i(r)t^{d-i},
 \qquad\Lambda_0=1,\quad\Lambda_d=D.
\]
The trace is \(\Lambda_1\), and the associated universal Cayley--Hamilton quotient enforces these characteristic equations. Naturality under scalar extension and multiplicativity, not merely a trace with formal divided powers, are essential.

The aid attributes this arbitrary-ring framework to Chenevier and exact agreement with Taylor trace pseudocharacters when \(d!\) is invertible to Ophir's theorem. The full references and hypotheses are not supplied. The framework is not a claim that every determinant over every ring is realized by an actual matrix representation over that same ring without further hypotheses. For group problems one applies it to \(R=A[G]\); continuity data require additional specification if \(G,A\) are topological.

\textbf{Reported degree-two trace fiber.} Fix a degree-two determinant \(D\) with trace \(T\). The aid reports that another determinant with the same trace is exactly \(D+q\), where \(q\) is a normalized additive homogeneous quadratic polynomial law with \(q(1)=0\), \(2q=0\), and
\[
 q(xy)=D(x)q(y)+D(y)q(x)+q(x)q(y).
\]
All these identities are required naturally after every commutative base change \(B\), for all \(x,y\in R\otimes_A B\); pointwise identities only on \(R\) are not substituted for polynomial-law identities. Conversely such a \(q\) gives a determinant with unchanged trace. This describes a possible nontrivial trace fiber in small characteristic rather than identifying trace and determinant data universally.

\textbf{Explicit reported-type characteristic-two obstruction.} Let \(A=\F_2[\epsilon]/(\epsilon^2)\) and \(C_4=\langle g:g^4=1\rangle\). One concrete realization uses
\[
 \rho_0(g)=I_2,\qquad
 \rho_1(g)=\begin{pmatrix}1&1\\\epsilon&1\end{pmatrix}.
\]
The second matrix squares to \((1+\epsilon)I_2\), hence has fourth power \(I_2\). Both representations have trace zero on every group element, and therefore identical trace laws on the group algebra, but
\[
 \det\rho_0(g)=1,\qquad\det\rho_1(g)=1+\epsilon.
\]
Their determinant polynomial laws are different. Thus even full trace data can lose infinitesimal deformation information when \(2!\) is not invertible. The example supports the need for the determinant-law replacement, not a failure of that replacement. All supplied framework, equivalence and fiber assertions require independent verification.
\subsection{Short English statement}
The aid reports determinant polynomial laws as the arbitrary-ring replacement for factorial-dependent trace pseudocharacters. In degree two a normalized quadratic torsion law can change the determinant without changing the trace. Two dual-number representations of \(C_4\) make that loss of deformation data explicit.
\subsection{Sources}
\begin{itemize}
\item Local UnsolvedMath registry, permanent ID \(20002836\), catalogue entry, \texttt{aim-representation-theory-notes.json}, \texttt{notes[107]}, workshop \emph{\(p\)-adic representations, modularity, and beyond}, Section 8 ``Kisin: Pseudo-representations'', Problem 1: \url{https://aimath.org/WWN/padicmodularity/padicmodularity.pdf}.
\item Supplied research attributions: Chenevier's multiplicative polynomial-law determinants and Cayley--Hamilton quotients; Ophir's agreement with Taylor pseudocharacters when \(d!\) is invertible. No full titles, dates or theorem numbers are supplied.
\item Supplied research summary on degree-two trace fibers and \(C_4\) dual-number representations. Attempt 1: full\_solution, status at run ``solved'', rigor ``proved'', low novelty confidence; registry solved/SOLVED-IN-LITERATURE. The machine-generated aid is not peer reviewed and requires independent verification before citation; no source visit is claimed.
\end{itemize}
\end{document}
